%%% Local Variables:
%%% mode: latex
%%% TeX-master: "../main"
%%% End:


\chapter{附录A}

%\section*{附录标题}
%\zihao{5}
%[内容为五号宋体。] 附录是作为论文主体的补充项目，并不是必须的。
%论文的附录依序用大写正体英文字母A、B、C……编序号，如：附录A。
%\vspace{0.75cm}
\begin{center}
\zihao{3}
\textbf{附录标题}
\end{center}



\indent
\zihao{5}
附录作为主体部分的补充，可包括以下内容：

——为了整篇论文材料的完整，但编入正文又有损于编排的条理性和逻辑性，这一材料包括比正文更为详尽的信息、研究方法和对技术更深入的叙述，对了解正文内容有用的补充信息等。

——由于篇幅过大或取材于复制品而不便于编入正文的材料。

——不便于编入正文的罕见珍贵资料。

——对一般读者并非必要阅读，但对本专业同行有参考价值的资料。

——正文中未被引用但被阅读或具有补充信息的文献。

——某些重要的原始数据、数学推导、结构图、统计表、计算机打印输出件等。


附录是作为论文主体的补充项目，并不是必须的。
论文的附录依序用大写正体英文字母A、B、C……编序号，如：附录A。